\section{Lateral Thinking Puzzles} \label{app:ltp}
\subsection{Dataset Details}
\vpara{Construction Details.}
Each sample is constructed of a pair of story (a riddle, e.g., \textit{A man walked into a restaurant, ordered a bowl of turtle soup, and after finishing it, he committed suicide. Why did he do that?}) and truth. 
We categorize samples into four levels of difficulty: easy, medium, hard, and expert. 
The LTP rules for LLM agent playing are as follows:

\begin{itemize}[leftmargin=1.5em,itemsep=0pt,parsep=0.2em,topsep=0.0em,partopsep=0.0em]
    \item \textbf{Roles:} Roles in LTP evaluation are a host and a solver. The host knows the story and truth, providing the story to the solver, and guiding it to guess out the truth. The solver, played and acted by an LLM, tries to find out the truth by asking questions and synthesizing host's answers.
    \item \textbf{Solving Steps:} There is a maximum round for each game, for example, 25. The solver needs to propose a question in each round based on known facts. The questions should be the ones that can be answered by ``Yes'', ``No'', or ``Irrelevant''. Host reply to the questions with correct answers. To lower the difficulty for LLM agents, sometimes the host will provides some hints in responses when solvers get trapped in wrong directions of reasoning.
    \item \textbf{Game Termination:} When the solver thinks it has guessed out the major part of the truth, it can declare the guessed plot to the host. If it is correct, the host will announce the end of the game.
\end{itemize}

\vpara{Evaluation Setup.}
For each pair of story and truth, we evaluate the models with the following steps:

\begin{itemize}[leftmargin=1.5em,itemsep=0pt,parsep=0.2em,topsep=0.0em,partopsep=0.0em]
    \item \textbf{Initialization.} Setting up the LTP host system via local python package installation or web API.
    
    \item \textbf{Interaction.} We set up system prompts for LLMs to build their roles of players. 
    LLMs are tested as solvers within the maximum round for each game, if the LLM does not exceed the max token length. 
    % For human evaluation, human provide feedback for questions asked by LLMs, and judge whether LLMs is asking effective questions and whether they have find out the truth. 
    In automatic evaluation, we limit the answer to be mostly "Yes", "No", or "Irrelevant", and extract the answer from \texttt{gpt-3.5-turbo}'s responses.
    LLMs are also asked to summarize their reasoning in automatic evaluation in order to help the termination detection to be more accurate.
    
    \item \textbf{Checking.} We do the pilot study of each LLM to collect all situations in game process and design the checking plan. For automatic evaluation, we set up some key words for \texttt{gpt-3.5-turbo} to answer and remind the model to consider some flexible situation like synonyms. 
    % We also double check the answer of gpt-3.5-turbo model to correct the bad answer. For human evaluation, we give each person games played by different models, to eliminate the bias between different people.
\end{itemize}

\vpara{Metrics.}
We evaluate LLMs' Lateral reasoning ability by two self created metrics:
\begin{itemize}[leftmargin=1.5em,itemsep=0pt,parsep=0.2em,topsep=0.0em,partopsep=0.0em]
    \item \textbf{Single Game Accuracy (SGA):} The proportion of rounds in which LLMs approaching the truth in a single game.
    \item \textbf{Round Efficiency (RE):} How fast the model can guess out the truth within the maximum round.
    \item \textbf{Query Relevance (QR):} Relevance between model's questions and the truth.
    \item \textbf{Game Progress (GP):} Progress before a game end, which serves as the main metric. We break down the groundtruth into several points and measure how many points are reached by an agent.
\end{itemize}

\subsection{Evaluation on LTP System}
We evaluate the LTP System by human validation, validating system's accuracy on milestone recognition and fact verification. We compare the Single Game Accuracy and Query Relevance between automatic evaluation and human evaluation, and found that automatic evaluation sometimes more tolerate for the agent, which make SGA and QR seem better than human evaluation, especially on open-sourced models. We plan to train a model specifically for the host of the game, in order to provide a better game experience and a more precise evaluation. For Game Progress and Round Efficiency, the LTP system provides an objective evaluation, which can match the level of human evaluation.

\subsection{LTP Game Progress and Termination}
The progress of game is defined as the proportion of hit key points in the truth. The key points are summarized by \texttt{gpt-3.5-turbo}, which are concluded in the dataset as ``answer\_keys'' (see an example below)

\lstset{
    backgroundcolor=\color[RGB]{245,245,244},
    breaklines=true,
    basicstyle=\ttfamily\small
}\begin{lstlisting}
Truth:
That night they went to the abandoned building to record the number of steps. They verified what was said on the Internet, and there would be one step less when counting the stairs at night. However, when they went to the abandoned building for verification the next day, they found that there were no stairs at all.}'':

Key points:
1. They want to count the steps of the abandoned building. 
2. A supernatural event occurred. 
3. They saw a claim online: counting stairs at night will result in one step less. 
4. The next day, when they went to the abandoned building to verify, they found no stairs. 
5. They broke down because they were terrified.

\end{lstlisting}
The number of key points varies among samples. As for the decision of whether the agent guess out key points, we first change relevant questions into declarative sentences, then simplify sentences into one sentence. After guessing out a key point, we delete that key point and relevant inferences to avoid repeated guessing. 
% An early termination of game is defined as guessing out more than 50\% of key points. 

\subsection{Prompt Example}
We use the following format of prompts for agents:
\lstset{
    backgroundcolor=\color[RGB]{245,245,244},
    breaklines=true,
    basicstyle=\ttfamily\small
}\begin{lstlisting}
You are a game player, and you are playing Lateral Thinking Puzzle, also known as Situation Puzzle.
Lateral Thinking Puzzle is a deductive reasoning game, and here are the game rules:
1. At the beginning of the game, you will receive a narrative, referred to as "story". Based on the story, you need to ask questions that can be answered with "yes", "no", or "irrelevant" to guees out the "truth".
2. By asking questions, you narrow down the range of possibilities until you eventually guess out the truth.
3. Each time, you can only ask one question.
4. Remember that your role is a player. You cannot declare the end of the game, give up on reasoning, or request a new game.
5. You cannot directly repeat information already provided in the story.
6. You cannot directly ask for details about the story in the form of "why" questions; you need to make your own guesses for truth.
7. You cannot directly inquire about the story; you must make your own deductions.

Next, please make full use of the information provided above to engage in game reasoning. Keep in mind that your questions should be answerable with "yes", "no", or "irrelevant", and you can only ask one question at a time.
Here is your story:
{story}

You can start guessing the content of the truth, and I will answer your questions. Please note that your questions should be answerable with "yes", "no", or "irrelevant".
\end{lstlisting}

We use the following format of prompts for host:
\lstset{
    backgroundcolor=\color[RGB]{245,245,244},
    breaklines=true,
    basicstyle=\ttfamily\small
}\begin{lstlisting}
USER:
I need you to be the host of a game called Lateral Thinking Puzzle.

Lateral Thinking Puzzle is a game consist of a story and a truth. Your story is: '{story}'
Your truth is: '{answer}'

Here are the game rules:
1. You know both the "story" and the "truth". When a user wants to play Lateral Thinking Puzzle, you provide them with the "story". The user only knows the "story" and is unawared of the "truth".
2. The user asks questions that can be answered with "yes," "no," or "irrelevant". Their questions are aimed at guessing the "truth". Based on the "truth", you respond to the user's questions using "yes," "no," or "irrelevant" to guide them towards guessing the correct truth.
3. If the user directly asks for details about the truth using the form of "why" questions, inform them that they need to make their own guesses.
4. You must fully understand and accurately interpret the information from the truth. Based on the information of the truth and the user's past questions, you answer the user's questions. The user's questions may not necessarily contain information from the truth, but your responses must align with the facts of the truth.
5. You can only answer "irrelevant" when the truth cannot provide a direct or indirect answer. Note that this is the only condition for responding "irrelevant"; otherwise, you should answer "yes" or "no."
6. You cannot directly disclose the information from the truth to the user, even if they ask directly.
7. You need to judge the user's questions as a whole and understand their overall intent. Avoid answering based solely on a particular point; your responses must align with the facts of the truth.
8. During the user's process of guessing the truth, if they come close to some truths but still have gaps in understanding the complete truth of the truth, you can provide certain entry point hints. However, you cannot directly reveal information from the truth.

During the game process, please adhere to the above game rules to ensure a positive gaming experience for the users. Pay close attention to the questions asked and ensure that your responses comply with both the game rules and the information from the truth. When a user requests to play the game, provide them with the story and help them guess the truth by answering with "yes", "no", or "irrelevant". Remember that with each response, you must fully understand and abide by the aforementioned game rules, as well as the story and the truth. This will ensure a smooth user experience and avoid situations where you cannot answer or violate the game rules.

ASSISTANT:
Alright, I understand that my role is to be the host of the Lateral Thinking Puzzle and help users guess the truth by answering their questions. I have fully grasped all the information regarding the story and the truth and have carefully read all the rules. I assure that I will abide by all the rules throughout the game process.

USER:
Please summarize the key points of the story to ensure that you have understood it.

ASSISTANT:
{story_key}

USER:
Please summarize the key points of the truth to ensure that you have understood it.

ASSISTANT:
{answer_key}

USER:
Please restate the rules to ensure that you have understood all of them.

ASSISTANT:
1. You know both the "story" and the "truth". When a user wants to play Lateral Thinking Puzzle, you provide them with the "story". The user only knows the "story" and is unawared of the "truth".
2. The user asks questions that can be answered with "yes," "no," or "irrelevant". Their questions are aimed at guessing the "truth". Based on the "truth", you respond to the user's questions using "yes," "no," or "irrelevant" to guide them towards guessing the correct truth.
3. If the user directly asks for details about the truth using the form of "why" questions, inform them that they need to make their own guesses.
4. You must fully understand and accurately interpret the information from the truth. Based on the information of the truth and the user's past questions, you answer the user's questions. The user's questions may not necessarily contain information from the truth, but your responses must align with the facts of the truth.
5. You can only answer "irrelevant" when the truth cannot provide a direct or indirect answer. Note that this is the only condition for responding "irrelevant"; otherwise, you should answer "yes" or "no."
6. You cannot directly disclose the information from the truth to the user, even if they ask directly.
7. You need to judge the user's questions as a whole and understand their overall intent. Avoid answering based solely on a particular point; your responses must align with the facts of the truth.
8. During the user's process of guessing the truth, if they come close to some truths but still have gaps in understanding the complete truth of the truth, you can provide certain entry point hints. However, you cannot directly reveal information from the truth.

USER:
Alright, we can now start the game. Remember, before each response, you should review the key points of the story, the key points of the truth, and the rules. Answer with "yes", "no", or "irrelevant".

ASSISTANT:
Alright, as the host of the game, I will adhere to the above rules and ensure that my responses comply with the rules and the information from the truth. Below is your story: 
{story}

You can start guessing the content of the truth, and I will answer your questions. Please note that your questions should be answerable with "yes", "no", or "irrelevant".

USER:
{question}
Please answer with "yes", "no", or "irrelevant".
\end{lstlisting}

Here is the prompt to convert questions answered by ``Yes'' into declarative sentence.
\lstset{
    backgroundcolor=\color[RGB]{245,245,244},
    breaklines=true,
    basicstyle=\ttfamily\small
}\begin{lstlisting}
Please restate the following content as a declarative sentence and simplify it into one sentence:
{question}
\end{lstlisting}

Here is the prompt to convert questions answered by ``No'' into declarative sentence.
\lstset{
    backgroundcolor=\color[RGB]{245,245,244},
    breaklines=true,
    basicstyle=\ttfamily\small
}\begin{lstlisting}
Please restate the following content as a declarative sentence by using the opposite meaning and then simplify it into one sentence:
{question}
\end{lstlisting}

Here is the prompt to merge reasoned out information into one sentence to judge whether the agent guess out the key points:
\lstset{
    backgroundcolor=\color[RGB]{245,245,244},
    breaklines=true,
    basicstyle=\ttfamily\small
}\begin{lstlisting}
Please simplify the following content into one sentence:
{reasoning}
\end{lstlisting}

Here is the prompt to judge whether the merged sentence hit the key point.
\lstset{
    backgroundcolor=\color[RGB]{245,245,244},
    breaklines=true,
    basicstyle=\ttfamily\small
}\begin{lstlisting}
Please compare the information between Sentence 1 and Sentence 2 to determine if Sentence 2 contains all the information in Sentence 1, including key details and descriptions. Please answer with "yes" or "no".
Sentence 1: {key}
Sentence 2: {merged sentence}"}
\end{lstlisting}